% Internal links for link previews: a contents page full of links (no previews there),
% then references to equations, a figure and footnotes across pages.
\documentclass[11pt,a4paper]{article}
\usepackage{amsmath}
\usepackage[hidelinks]{hyperref}
\usepackage{lipsum}
\begin{document}
\tableofcontents
\newpage
\section{Laplace transform}
The Laplace transform of $f$ is defined in \eqref{eq:laplace}; its linearity follows from
\eqref{eq:linear}, and the final value theorem \eqref{eq:final} closes the section.\footnote{See also the table of transforms in Figure~\ref{fig:table}.}
Stability is discussed in Section~\ref{sec:stab}.
\lipsum[1-2]
\newpage
\section{Definitions}
\begin{equation}\label{eq:laplace}
  F(s) = \int_0^\infty f(t)\,e^{-st}\,dt
\end{equation}
\lipsum[3]
\begin{equation}\label{eq:linear}
  \mathcal{L}\{a f + b g\} = a F(s) + b G(s)
\end{equation}
\lipsum[4]
\begin{equation}\label{eq:final}
  \lim_{t\to\infty} f(t) = \lim_{s\to 0} s F(s)
\end{equation}
\newpage
\section{Stability}\label{sec:stab}
\lipsum[5]
\begin{figure}[h]
  \centering
  \fbox{\parbox{0.6\textwidth}{\centering Table of common transforms\\[4pt] $1 \leftrightarrow 1/s$ \quad $e^{at} \leftrightarrow 1/(s-a)$}}
  \caption{Common Laplace transforms}\label{fig:table}
\end{figure}
\lipsum[6]

\section{Poles} \lipsum[7]
\section{Zeros} \lipsum[8]
\section{Bode plots} \lipsum[9]
\section{Nyquist} \lipsum[10]
\section{Root locus} \lipsum[11]
\section{PID control} \lipsum[12]
\section{State space} \lipsum[13]
\section{Observers} \lipsum[14]
\section{Discrete time} \lipsum[15]
\section{Z transform} \lipsum[16]
\section{Sampling} \lipsum[17]
\section{Exercises} \lipsum[18]
\end{document}
